\documentclass[11pt]{article}

% === Packages ===
\usepackage[utf8]{inputenc}
\usepackage[T1]{fontenc}
\usepackage{times}
\usepackage[margin=1in]{geometry}
\usepackage{amsmath,amssymb}
\usepackage{graphicx}
\usepackage{booktabs}
\usepackage{multirow}
\usepackage{hyperref}
\usepackage{xcolor}
\usepackage{enumitem}
\usepackage{caption}
\usepackage{subcaption}
\usepackage{natbib}
\usepackage{algorithm}
\usepackage{algpseudocode}

\hypersetup{colorlinks=true,linkcolor=blue,citecolor=blue,urlcolor=blue}

% === Title ===
\title{Naive Stepwise Partial Rewards Induce Early-Stop Collapse\\in GRPO Training of Sequential API Agents}

\author{
  [Author Name]\\
  [Affiliation]\\
  \texttt{[email]}
}

\date{}

\begin{document}
\maketitle

% ============================================================
\begin{abstract}
We study reward design for GRPO training of sequential API agents.
Using a verifiable order management gym with deterministic state transitions, we train Qwen2.5-0.5B-Instruct across seven reward variants and evaluate with multi-seed protocols (up to 900 tasks per model).
Our main finding is a negative result: naive stepwise partial reward---which grants per-step credit based on format compliance without API verification---achieves only $11.0\% \pm 1.0\%$ overall success, compared with $73.6\% \pm 1.0\%$ for an outcome-dominant reward that conditions the maximum signal on verified task completion.
The naive partial model exhibits early-stop collapse: on hard tasks (4--5 steps), it stops after $\sim$1 step in 100\% of trials across all seeds, collecting partial credit rather than attempting completion.
Note that this comparison simultaneously changes two variables---reward verification and reward structure---so the result should be understood as evidence against \emph{naive unverified stepwise} partial reward specifically, not against all forms of dense or intermediate reward.
We additionally find that initializing from an outcome-dominant checkpoint reduces sensitivity to reward type in subsequent training, with continuation models converging to $\sim$82\% regardless of reward signal.
All experiments run on a single consumer laptop GPU (RTX~5060, 8\,GB VRAM).
\end{abstract}

% ============================================================
\section{Introduction}
\label{sec:intro}

Reinforcement Learning from Verifiable Rewards (RLVR) has emerged as a powerful paradigm for training large language models, with systems like DeepSeek-R1 \citep{deepseek2025r1} demonstrating that binary correctness signals alone can induce sophisticated reasoning behaviors.
However, nearly all RLVR research focuses on mathematical \citep{shao2024deepseekmath,hendrycks2021math,cobbe2021gsm8k} and coding tasks with single-turn verification.
Real-world deployment increasingly requires \emph{agents} that execute multi-step workflows---calling APIs, managing state transitions, and coordinating across entities \citep{yao2024taubench,agentr1,agentrl2025}---where the gap between ``partial progress'' and ``task completion'' is both natural and tempting to exploit as a training signal \citep{casper2023rewardhacking}.

The intuition is compelling: if a task requires five sequential API calls, surely rewarding a model that completes three steps should provide more training signal than a binary pass/fail.
Step-level feedback should help with credit assignment, reduce reward sparsity, and guide the policy toward correct behavior.
This intuition has driven process reward models \citep{lightman2023lets,wang2024mathshepherd}, stepwise verification \citep{uesato2022solving}, and turn-level reward design for multi-turn RL \citep{zeng2025turnlevel}.

We test this intuition in a controlled setting and find that \emph{naive} stepwise partial credit---based on format compliance without API verification---can be severely counterproductive, consistent with findings in mathematical reasoning \citep{gao2024designing} but, to our knowledge, not previously documented for sequential API agent tasks.

\paragraph{Research Question.}
In deterministic API workflows with strict order dependencies, does naive stepwise partial credit improve GRPO training, or does outcome-dominant reward suffice?

\paragraph{Setup.}
We build a \emph{Zendesk-style training gym}: a verifiable order management API with deterministic state machine transitions, a parametric task generator spanning three difficulty levels (1--5 steps), and an automatic verifier that checks final order states.
We train Qwen2.5-0.5B-Instruct \citep{qwen2025qwen25} with TRL's GRPOTrainer \citep{trl2024}, building on the GRPO algorithm introduced for DeepSeek \citep{deepseek2025r1}, across seven variants: format-only baseline, outcome-dominant reward, outcome-dominant with Boltzmann curriculum, naive step-level partial credit, stepwise execution with partial credit, supervised fine-tuning, and a verifiable milestone hybrid.
All models are evaluated with a unified 30-task protocol across three random seeds.

\paragraph{Findings.}
On expanded evaluation (100 tasks per difficulty, 3 seeds = 900 tasks per model), outcome-dominant reward achieves $73.6\% \pm 1.0\%$ overall versus $11.0\% \pm 1.0\%$ for naive stepwise partial.
On hard tasks (4--5 steps), the gap is absolute: $31.3\%$ vs. $0.0\%$.
The partial-reward model exhibits early-stop collapse \citep{skalse2022misspecification}: a 100\% early-stop failure rate on hard tasks across all three seeds, completing an average of 1.07 out of 4.3 required steps before stopping.
Note that this comparison confounds two changes---reward verification and reward structure (Section~\ref{sec:head2head})---so we interpret this as a negative result about naive unverified stepwise partial reward, not a general claim against dense rewards.
Continuation training from an outcome-dominant checkpoint converges to $\sim$82\% regardless of reward type, suggesting that strong initialization reduces sensitivity to reward signal.

\paragraph{Contributions.}
\begin{enumerate}[leftmargin=*,itemsep=1pt]
  \item A \textbf{verifiable API workflow gym} for RLVR agent training: a deterministic order management state machine with parametric task generation, automatic verification, and three difficulty levels.
  \item An \textbf{early-stop collapse finding}: in our setting, naive stepwise partial reward induces a stereotyped policy that stops after $\sim$1 step on multi-step tasks, with 100\% reproducibility across seeds. We introduce a failure taxonomy (Type~A suboptimal convergence vs. Type~B reward-induced stopping) to characterize this.
  \item An \textbf{initialization effect}: models initialized from outcome-dominant checkpoints are substantially less sensitive to reward type during continuation training.
\end{enumerate}

% ============================================================
\section{Related Work}
\label{sec:related}

\paragraph{RLVR and GRPO.}
DeepSeek-R1 \citep{deepseek2025r1} demonstrated that RLVR with binary rewards can train reasoning capabilities without supervised data.
\citet{mroueh2025grpo} provided theoretical analysis of GRPO under binary rewards, while \citet{wen2025rlvr} showed RLVR incentivizes correct reasoning chains and \citet{yue2025rlvr} questioned whether RLVR's impact persists beyond training distribution.
\citet{huang2026dynamics} developed a theory of RLVR training dynamics showing that smooth difficulty spectra produce a ``relay effect'' enabling steady improvement, while abrupt difficulty discontinuities cause grokking-type plateaus.
\citet{liu2025drgrpo} identified an optimization bias in GRPO that artificially inflates response length for incorrect outputs, proposing Dr. GRPO as an unbiased fix.
DAPO \citep{yu2025dapo} documented entropy collapse in large-scale GRPO training and introduced decoupled clipping and dynamic sampling as mitigations.
For a comprehensive survey of reward mechanisms and their failure modes in RL-based reasoning, see \citet{pan2026survey}.

\paragraph{Process vs. Outcome Rewards.}
\citet{gao2024designing} showed that process rewards can \emph{hurt} RL training when combined with success rewards in PPO \citep{schulman2017ppo}.
\citet{sullivan2025grpoprm} proved that GRPO implicitly functions as a process reward model, suggesting that explicit step-level signals may double-count credit assignment.
PRIME \citep{cui2025prime} demonstrated implicit process rewards from outcome labels. UltraLogic \citep{ultralogic2025} introduced bipolar graded float rewards (positive for correct, negative for incorrect reasoning) and reported gains over standard 0/1 binary---a finding that supports richer reward signals when properly structured, contrasting with our negative result for naive unverified stepwise credit.
MO-GRPO \citep{mogrpo2025} identified multi-objective reward hacking where GRPO optimizes the highest-variance reward component.
\citet{chen2026milestones} replaced sparse outcome signals with dense sub-goal verifiable rewards for geometric reasoning, achieving strong gains---but crucially, each sub-goal remains individually verifiable, unlike naive stepwise partial credit.

\paragraph{Reward Design for Agent RL and Tool-Use.}
Recent work has shown that reward design is central to RL for long-horizon language agents and tool-use systems, but the emerging consensus is nuanced rather than binary.
ToolRL \citep{toolrl2025} studies reward design for tool-using language agents and shows that reward type, scale, granularity, and temporal dynamics all matter for stable learning.
CM2 \citep{cm2_2025} further argues that naively densifying rewards in multi-turn, multi-step agentic tool use can amplify noise and destabilize optimization, motivating checklist-style structured credit assignment instead.
In adjacent settings, LongRLVR \citep{zhang2025longrlvr} introduces verifiable context rewards for long-context RL, ADMIRE \citep{admire2025} proposes adaptive milestone rewards for GUI agents, and Search-P1 \citep{searchp1} uses path-centric reward shaping for agentic retrieval.
SETA \citep{li2025seta} uses independent unit test pass ratio with a completion bonus for software engineering agents---their tests are independently verifiable (unlike our sequential steps), making early stopping non-exploitable.
\citet{seta2026} uses unit test pass ratio as dense reward for terminal agents trained with GRPO, explicitly rejecting binary 0/1 as ``too sparse for calculating advantage,'' while Endless Terminals \citep{gandhi2026endless} achieves strong gains with purely binary episode-level rewards by scaling environment diversity to 3,255 tasks---supporting our finding that binary rewards can suffice when environments are sufficiently diverse.
WebAgent-R1 \citep{wei2025webagent} extends GRPO to multi-turn web agent training with binary rewards, boosting Qwen-2.5-3B from 6.1\% to 33.9\% on WebArena-Lite.
Planner-R1 \citep{zhu2025planner} finds that smaller models are highly responsive to dense reward shaping in agentic RL, achieving 56.9\% pass rate on TravelPlanner with only 180 training queries.
\citet{su2025prs} proposes Progressive Reward Shaping (PRS), a curriculum-inspired dense stage-wise reward, and VSPO, an enhanced GRPO variant that replaces zero-advantage samples.
iStar \citep{liu2025istar} introduces implicit step rewards via a trajectory-based DPO objective, noting that explicit process supervision in agent RL suffers from biased annotation and reward hacking.
Agent-R1 \citep{agentr1}, AgentRL \citep{agentrl2025}, $\tau$-bench \citep{yao2024taubench}, and WebRL \citep{qi2024webrl} further study multi-turn agent RL and evaluation.
Taken together, these works suggest that dense intermediate feedback \emph{can} help, but only when it is structured, verifiable, and aligned with final task completion---rather than provided as free-form stepwise partial credit.
The emerging landscape reveals a striking divergence: \citet{seta2026} and \citet{zhu2025planner} favor dense rewards, while \citet{gandhi2026endless}, \citet{wei2025webagent}, and our work find binary rewards sufficient or superior---suggesting that the optimal reward granularity depends critically on task structure and step independence.

\paragraph{Positioning.}
Our work focuses on a narrower but practically important setting: deterministic, sequential API workflows with strict order dependencies and verifiable end states.
Unlike prior work that studies reward design broadly for tool use or introduces structured dense rewards, we isolate a specific failure mode of naive stepwise partial reward under training from scratch: \emph{early-stop collapse}.
In our benchmark, the model trained with naive partial credit collapses to trajectories that terminate after roughly one step on hard tasks, yielding 0\% D2/D3 success and 100\% Type~B early-stop failures on D3, while an outcome-dominant reward remains stable.
We further show that this vulnerability is strongly initialization-dependent: continuation from a strong checkpoint substantially reduces reward sensitivity.
To our knowledge, this combination of findings---early-stop collapse in deterministic sequential API workflows, quantified through failure taxonomy and average completed steps, together with a from-scratch vs. continuation contrast---has not been directly documented in prior work.


% ============================================================
\section{The Zendesk Training Gym}
\label{sec:gym}

We construct a verifiable training environment for API agent tasks, inspired by real customer service workflows such as those modeled by $\tau$-bench \citep{yao2024taubench}. Our gym consists of three components: an API system with deterministic state transitions, a parametric task generator, and an automatic verifier.

\subsection{API System}

The core is an order management REST API implementing a finite state machine:
\begin{center}
\texttt{created} $\to$ \texttt{paid} $\to$ \texttt{shipped} $\to$ \texttt{delivered} $\to$ \texttt{refunded}\\[2pt]
\texttt{created} $\to$ \texttt{cancelled}
\end{center}

Each transition requires specific parameters (e.g., payment method, carrier, tracking number) and is enforced deterministically---invalid transitions return structured error responses.
The API is implemented with FastAPI and deployed on localhost during training, ensuring zero network latency and full reproducibility.

\subsection{Task Generator}

Tasks span three difficulty levels with nine task types:
\begin{itemize}[leftmargin=*,itemsep=1pt]
  \item \textbf{D1 (Easy, 1--2 steps):} \texttt{create\_order}, \texttt{create\_and\_check}, \texttt{create\_and\_cancel}
  \item \textbf{D2 (Medium, 2--3 steps):} \texttt{create\_and\_pay}, \texttt{create\_pay\_ship}, \texttt{pay\_and\_refund}
  \item \textbf{D3 (Hard, 4--5 steps):} \texttt{full\_delivery}, \texttt{multi\_order}, \texttt{full\_lifecycle\_with\_refund}
\end{itemize}

Parameters (customer names, products, prices, carriers) are randomized per task.
Tasks are generated procedurally rather than drawn from a fixed pool: training and evaluation instances are sampled independently from the generator using separate random draws, so the evaluation set contains no exact duplicate of any training instance.
The train and evaluation distributions share the same task schema and API semantics but differ at the level of instantiated parameters.
The \texttt{multi\_order} type requires managing two parallel orders with different lifecycle targets, testing entity tracking capability.

\subsection{Verifier}

The verifier: (1)~parses the model's output for structured \texttt{ACTION}/\texttt{BODY} pairs; (2)~executes parsed actions against the live API; (3)~checks that the final order state matches the task specification.
It returns both a binary success signal and a partial reward $r_{\text{partial}} = k / T$, where $k$ is the number of completed steps and $T$ is the total required, enabling comparison between reward types using the same underlying evaluation.

\subsection{Boltzmann Curriculum}
\label{sec:curriculum}

For variant v2b, we implement adaptive difficulty sampling.
At each training step, difficulty $d$ is sampled with probability proportional to $\exp(-\text{perf}_d / \tau)$, where $\text{perf}_d$ is the running success rate for difficulty $d$ and temperature $\tau$ anneals from 2.0 to 0.3 over training, progressively focusing on the hardest tasks.


% ============================================================
\section{Experimental Setup}
\label{sec:setup}

\subsection{Model and Training}

We use Qwen2.5-0.5B-Instruct \citep{qwen2025qwen25} as the base model, selected for its chat-format pre-training and feasibility on consumer hardware.
Training uses TRL's GRPOTrainer \citep{trl2024}, which implements the GRPO algorithm \citep{shao2024deepseekmath,deepseek2025r1}, with: bfloat16 precision, gradient checkpointing, batch size 1 with gradient accumulation over 4 steps, 100 tasks per epoch, 3 epochs, 4 generations per prompt, and KL coefficient $\beta = 0.04$.
All training runs are conducted on a single NVIDIA RTX~5060 Laptop GPU with 8\,GB VRAM, with most variants completing in 60--90 minutes and v5 (milestone hybrid) requiring approximately 3 hours due to its additional verification overhead.

\subsection{Training Variants}

Table~\ref{tab:variants} describes our seven training variants.
The key comparison is between variants trained \emph{from scratch} (from the base Qwen2.5 checkpoint) and \emph{continuation} variants (initialized from a previously trained model).

\begin{table}[t]
\centering
\caption{Training variants. ``Scratch'' = trained from Qwen2.5-0.5B-Instruct; ``Cont.'' = continuation from an outcome-dominant-trained checkpoint.}
\label{tab:variants}
\small
\begin{tabular}{@{}llll@{}}
\toprule
\textbf{Variant} & \textbf{Reward} & \textbf{Init} & \textbf{Key Design} \\
\midrule
v1   & format only          & Scratch & Baseline GRPO \\
v2a  & format + outcome-dominant & Scratch & Success$\to$1.0, fail$\to$partial \\
v2b  & format + outcome-dom. + curric. & Cont. (v2a) & + Boltzmann curriculum \\
v2c  & format + naive step partial & Scratch & Per-step format credit \\
v3a  & SFT                  & Scratch & Supervised fine-tuning \\
v4   & format + stepwise partial & Cont. (v2b) & Step-by-step + partial \\
v5   & format + milestone hybrid & Scratch & $0.7 r_{\text{bin}} + 0.3 (m/M)$ \\
\bottomrule
\end{tabular}
\end{table}

The \textbf{outcome-dominant reward} (v2a) is a hybrid: it returns 1.0 on success (final state matches specification) and $r_{\text{partial}} = k/T$ on failure (where $k$ is completed steps and $T$ is total required).
This is \emph{not} pure binary---a model completing 4/5 steps but failing final verification receives $r = 0.80$, not $0.0$.
The key property is \emph{outcome dominance}: the gap between ``almost done'' ($\sim$0.80) and ``done'' (1.0) is $\sim$0.20, which still incentivizes full completion.
The \textbf{naive stepwise partial reward} (v2c) returns per-step format credit without API verification: each step that follows the \texttt{ACTION}/\texttt{BODY} format receives equal credit regardless of whether it causes a correct state transition.
Both variants also include a format reward (0.0--1.0) that checks the overall output structure.

\subsection{Evaluation Protocol}

Each model is evaluated on 30 tasks (10 per difficulty level) using a unified evaluation script with deterministic task generation per seed.
All evaluation tasks are newly sampled instances not reused from training.
We report results across three random seeds (42, 123, 777) for the four key variants (v2a, v2b, v2c, v4), yielding 12 total evaluations.
Variants v1 and v3a use single-seed evaluation as reference baselines.

\paragraph{Failure Taxonomy.}
We classify each failure into two primary categories aligned with the reward hacking literature \citep{casper2023rewardhacking,singhal2023rewardhacking}:
\begin{itemize}[leftmargin=*,itemsep=1pt]
  \item \textbf{Type~A} (``high partial, low final''): The model generates actions for most or all required steps ($r_{\text{partial}} \geq 0.75$) but fails final state verification. This indicates suboptimal policy convergence---the model \emph{tries} but makes parameter errors.
  \item \textbf{Type~B} (``early stop''): The model generates fewer actions than required. This indicates reward-hacking-induced early stopping---the model \emph{chooses not to try}, collecting partial credit instead.
\end{itemize}


% ============================================================
\section{Results}
\label{sec:results}

\subsection{Main Results}

Figure~\ref{fig:overall} and Table~\ref{tab:main} present multi-seed results for our four primary variants, alongside single-seed results for reference baselines.

\begin{figure}[t]
\centering
\includegraphics[width=0.85\linewidth]{fig1_overall.pdf}
\caption{Overall success across training variants (30-task eval, 3 seeds). Reward effectiveness is non-monotonic: neither naive stepwise partial credit (v2c) nor a simple verifiable-milestone hybrid (v5) matches the balance achieved by the outcome-dominant variant (v2a), suggesting the key issue is reward balance and executability rather than reward density per se. Error bars show $\pm$1 std across 3 seeds.}
\label{fig:overall}
\end{figure}

\begin{table}[t]
\centering
\caption{Main results. Multi-seed variants report mean $\pm$ std across 3 seeds. D1/D2/D3 = easy/medium/hard difficulty. $^\dagger$Single-seed only.}
\label{tab:main}
\small
\begin{tabular}{@{}llccccc@{}}
\toprule
\textbf{Variant} & \textbf{Reward} & \textbf{Init} & \textbf{D1} & \textbf{D2} & \textbf{D3} & \textbf{Overall} \\
\midrule
v1$^\dagger$ & format only & Scratch & 100 & 90 & 0 & 63 \\
v3a$^\dagger$ & SFT & Scratch & 100 & 40 & 0 & 46 \\
\midrule
\textbf{v2a} & \textbf{outcome-dom.} & \textbf{Scratch} & $90{\pm}8$ & $100{\pm}0$ & $\mathbf{47{\pm}5}$ & $\mathbf{79{\pm}3}$ \\
v2c & naive partial & Scratch & $43{\pm}9$ & $0{\pm}0$ & $0{\pm}0$ & $14{\pm}3$ \\
v5 & milestone hybrid & Scratch & $0{\pm}0$ & $0{\pm}0$ & $0{\pm}0$ & $0{\pm}0$ \\
\midrule
v2b & outcome-dom. + curric. & Cont. & $97{\pm}5$ & $100{\pm}0$ & $30{\pm}8$ & $76{\pm}4$ \\
v4 & stepwise partial & Cont. & $100{\pm}0$ & $100{\pm}0$ & $47{\pm}5$ & $82{\pm}2$ \\
\bottomrule
\end{tabular}
\end{table}

\paragraph{From-Scratch: Outcome-Dominant $\gg$ Naive Partial.}
Figure~\ref{fig:overall} compares all training variants under the 30-task evaluation protocol (10 tasks per difficulty, 3 seeds; see Section~\ref{sec:setup}).
Figure~\ref{fig:overall} shows a non-monotonic pattern rather than a simple ordering by ``reward sophistication'': both naive stepwise partial reward (v2c, $14.4\%$) and a simple verifiable-milestone hybrid (v5, $0.0\%$) underperform the outcome-dominant variant (v2a, $78.9\%$), suggesting that the key issue is reward balance and executability, not merely whether intermediate reward is present---echoing \citet{gao2024designing}'s finding that process rewards can hurt RL training even when they seem informative.
Note that v2c reports $14.4\%$ here under 30-task evaluation; the expanded 900-task evaluation in Section~\ref{sec:head2head} yields a tighter estimate of $11.0\% \pm 1.0\%$---the difference reflects the higher variance of the smaller eval set, not a discrepancy.

\paragraph{Continuation Models Converge.}
In contrast, continuation training from an outcome-dominant-trained checkpoint yields nearly identical results regardless of reward type: v4 (partial, continuation) achieves $82.2\% \pm 1.6\%$, comparable to v2a (outcome-dominant, from scratch) at $78.9\% \pm 3.1\%$.
The prior knowledge learned from outcome-dominant training---specifically, the behavior of attempting all required steps---persists through continuation training with partial reward, consistent with findings that strong initialization can stabilize RL fine-tuning against reward exploitation \citep{toolrl2025}.

\paragraph{RLVR Outperforms SFT.}
Even the format-only GRPO baseline (v1, 63\%) outperforms supervised fine-tuning (v3a, 46\%), confirming the RLVR advantage for agent tasks with minimal reward engineering---paralleling the reasoning-domain finding that RLVR can match or exceed supervised approaches \citep{deepseek2025r1,wen2025rlvr}.

\subsection{The Critical Comparison: v2a vs. v2c}
\label{sec:head2head}

To validate the core comparison with higher statistical resolution, we conduct an expanded evaluation with 100 tasks per difficulty level (300 total per seed, 3 seeds = 900 tasks per model).
Table~\ref{tab:head2head} presents the results.

\paragraph{Important caveat.}
The v2a vs. v2c comparison simultaneously changes two variables: (1)~reward \emph{verification}---v2a's partial-on-failure component checks API state, while v2c grants credit based on format compliance alone; and (2)~reward \emph{structure}---outcome-dominant vs. stepwise.
Consequently, we cannot attribute the performance gap to reward granularity alone.
Our result is best understood as a negative finding about \emph{naive unverified stepwise} partial reward, not as evidence that all forms of dense or intermediate reward are harmful.
Structured, verified dense rewards \citep{cm2_2025,admire2025,chen2026milestones} may well avoid this failure mode.

\begin{table}[t]
\centering
\caption{Head-to-head: outcome-dominant (v2a) vs. naive partial (v2c), both trained from scratch. Expanded evaluation: 100 tasks per difficulty $\times$ 3 seeds = 900 tasks per model.}
\label{tab:head2head}
\small
\begin{tabular}{@{}lrrr@{}}
\toprule
\textbf{Metric} & \textbf{v2a (outcome-dom.)} & \textbf{v2c (naive partial)} & \textbf{Gap} \\
\midrule
Overall success & $73.6{\pm}1.0\%$ & $11.0{\pm}1.0\%$ & $6.7\times$ \\
D1 success & $90.0{\pm}0.8\%$ & $33.0{\pm}2.9\%$ & $2.7\times$ \\
D2 success & $99.3{\pm}0.5\%$ & $0.0{\pm}0.0\%$ & $\infty$ \\
D3 success & $31.3{\pm}4.1\%$ & $0.0{\pm}0.0\%$ & $\infty$ \\
D3 steps completed & $3.47{\pm}0.03$ & $1.06{\pm}0.01$ & $3.3\times$ \\
Early stop rate (all) & $12.2{\pm}1.1\%$ & $89.0{\pm}1.0\%$ & $7.3\times$ \\
D3 Type~B rate & $36.7{\pm}3.4\%$ & $100.0{\pm}0.0\%$ & $2.7\times$ \\
\bottomrule
\end{tabular}
\end{table}

Figure~\ref{fig:expanded} strengthens the main finding under expanded evaluation. The gap between v2a and v2c widens from $5.5\times$ (30-task eval) to $6.7\times$ (300-task eval), while standard deviations shrink 3--5$\times$.
Crucially, the gap is not confined to the hardest tasks: v2c also collapses on D2, where it achieves $0.0\%$ success compared with $99.3\%$ for v2a.
This suggests that the failure is triggered not merely by task difficulty, but by sequential dependency structure itself: once tasks require ordered completion of interdependent API operations, naive partial credit becomes highly unstable during training from scratch \citep{cm2_2025}.
The partial-reward model completes an average of 1.06 steps on D3 tasks requiring 4--5 steps and achieves $100.0\%$ D3 Type~B rate with $\sigma = 0.0\%$ across all three seeds, ruling out sampling variance. This is consistent with reward hacking: the learned policy exploits the format-based partial credit signal rather than attempting full task completion \citep{skalse2022misspecification,casper2023rewardhacking}.

\begin{figure}[t]
\centering
\includegraphics[width=0.85\linewidth]{fig2_expanded.pdf}
\caption{Expanded evaluation comparing the outcome-dominant variant (v2a) with naive stepwise partial reward (v2c). 100 tasks per difficulty $\times$ 3 seeds = 900 tasks per model. The collapse of v2c extends beyond the hardest tasks and appears whenever ordered completion of dependent API operations is required.}
\label{fig:expanded}
\end{figure}

\subsection{Outcome-Dominant Reward Eliminates Early Stopping}

Variant v2b (outcome-dominant + Boltzmann curriculum) provides additional evidence.
Across all three seeds, v2b achieves \textbf{0\% Type~B (early stop) failures}---on any difficulty level.
All 21 D3 failures are Type~A (high partial, 86\%) or multi-entity management errors (14\%).
The model always \emph{attempts} every required step; it sometimes fails final verification due to parameter errors, but never stops early.

This directly contrasts with v2c's 100\% Type~B rate: outcome-dominant reward produces an ``attempt everything'' policy, while naive partial reward produces a ``harvest partial credit'' policy---consistent with \citet{sullivan2025grpoprm}'s argument that GRPO's implicit process reward mechanism may be sufficient for credit assignment without explicit step-level signals.

% ============================================================
\section{Analysis: Reward Hacking Mechanism}
\label{sec:analysis}

\subsection{Multi-Seed Failure Taxonomy}

Table~\ref{tab:taxonomy} presents the aggregated D3 failure taxonomy across all seeds.

\begin{table}[t]
\centering
\caption{D3 failure taxonomy aggregated across 3 seeds. Total D3 tasks = 30 per model (10 tasks $\times$ 3 seeds). ``Other'' for v5 = semantic failure (syntactically valid actions, 0 verified steps).}
\label{tab:taxonomy}
\small
\begin{tabular}{@{}lcccc@{}}
\toprule
\textbf{Model} & \textbf{D3 Fails} & \textbf{Type~A} & \textbf{Type~B} & \textbf{Other} \\
\midrule
v2a (outcome-dom., scratch) & 16 & 8 (50\%) & 8 (50\%) & 0 \\
v2b (outcome-dom., cont.) & 21 & 18 (86\%) & \textbf{0 (0\%)} & 3 (14\%) \\
v2c (naive partial, scratch) & \textbf{30} & 0 (0\%) & \textbf{30 (100\%)} & 0 \\
v5 (milestone hyb., scratch) & \textbf{30} & 0 (0\%) & 10 (33\%) & \textbf{20 (67\%)} \\
v4 (partial, cont.) & 16 & 8 (50\%) & 8 (50\%) & 0 \\
\bottomrule
\end{tabular}
\end{table}

Figure~\ref{fig:taxonomy} visualizes the taxonomy and reveals that the difference is not only in success rate, but in \emph{failure regime}---reward design changes the shape of the learned policy, not just its final accuracy. Three patterns emerge:

\begin{figure}[t]
\centering
\includegraphics[width=0.85\linewidth]{fig3_taxonomy.pdf}
\caption{D3 failure taxonomy by model. Reward design changes not only success rate but also failure regime: v2c collapses entirely to Type~B early stopping, v2b eliminates Type~B failures, v2a/v4 show a balanced Type~A/B split, and v5 exhibits a distinct pathology dominated by semantic failure (syntactically valid actions that fail API verification, yielding 0 verified steps). 30 D3 tasks = 10 $\times$ 3 seeds.}
\label{fig:taxonomy}
\end{figure}

\paragraph{v2c: Pure Reward Hacking.}
All 30 D3 evaluations (10 tasks $\times$ 3 seeds) fail, and every single failure is Type~B (early stop).
The model learns a stereotyped policy: execute one action and stop, exploiting the format-based partial credit signal \citep{skalse2022misspecification}, analogous to the length-exploitation phenomenon documented in RLHF settings \citep{singhal2023rewardhacking}.
With D3 tasks requiring 4--5 steps, the average completion is 1.07 steps---the model barely begins the task before stopping.
This is not variance or bad luck; the standard deviation across seeds is 0.0\%.

\paragraph{v2b: Zero Early Stopping.}
Binary reward with curriculum produces 0\% Type~B failures across all seeds.
All failures are Type~A (the model attempts all steps but makes errors) or multi-entity tracking failures.
Outcome-dominant reward trains the model to \emph{always attempt completion}, even at the cost of sometimes failing verification \citep{deepseek2025r1}.

\paragraph{v5: A Different Pathology.}
V5 also fails all 30 D3 tasks, but its failure distribution differs sharply from v2c: 67\% are semantic failures (the model generates syntactically well-formed actions whose API calls consistently fail verification, yielding 0 verified steps) and only 33\% are Type~B early stops.
Where v2c learns to \emph{stop trying}, v5 learns to \emph{try incorrectly}---optimizing surface-form validity without acquiring grounded API semantics \citep{mogrpo2025,fontana2025hidden}.

\paragraph{v4 $\approx$ v2a: Strong Initialization.}
Despite receiving partial reward during training, v4 exhibits the same failure distribution as v2a (50/50 Type~A/B split).
The outcome-dominant-trained checkpoint's behavioral prior---``attempt all steps''---reduces sensitivity to the partial reward signal during continuation training.

\subsection{Characteristic Stopping Patterns}

The v2c model's early stopping is highly stereotyped:
\begin{itemize}[leftmargin=*,itemsep=1pt]
  \item \texttt{full\_lifecycle\_refund} (5 steps): completes 1 step $\to$ partial $\approx 0.20$
  \item \texttt{multi\_order} (4 steps): completes 1 step $\to$ partial $\approx 0.25$
  \item \texttt{full\_delivery} (4 steps): completes 1 step $\to$ partial $\approx 0.25$
\end{itemize}

For comparison, the first v4 run (partial, from checkpoint) showed different stopping points:
\begin{itemize}[leftmargin=*,itemsep=1pt]
  \item \texttt{full\_lifecycle\_refund} (5 steps): stops at 3 $\to$ partial $= 0.60$
  \item \texttt{multi\_order} (4 steps): stops at 2 $\to$ partial $= 0.50$
  \item \texttt{full\_delivery} (4 steps): stops at 3 $\to$ partial $= 0.75$
\end{itemize}

The from-scratch model (v2c) stops earlier because it never learned multi-step execution.
The continuation model (first v4 run) stops later---it inherited some capability from outcome-dominant training---but still converges to a partial-credit optimum rather than attempting full completion, illustrating how partial reward can gradually erode learned behaviors even from a strong starting point \citep{gao2024designing}.

\subsection{Connection to Theory}

Our empirical findings align with several theoretical results:

\paragraph{GRPO as Implicit PRM \citep{sullivan2025grpoprm}.}
If GRPO already functions as an implicit process reward model, then adding explicit step-level rewards creates double-counting.
This explains why v2a (outcome-dominant) outperforms v2c (naive stepwise partial): the explicit per-step signal \emph{interferes} with GRPO's internal credit assignment.

\paragraph{MO-GRPO Reward Hacking \citep{mogrpo2025}.}
When GRPO optimizes multiple reward components, it preferentially optimizes the highest-variance component.
In our setup, the step partial reward has higher variance (0.0--1.0 in fine increments) than the binary reward (0 or 1), potentially causing the optimizer to focus on maximizing partial credit at the expense of binary success.

\paragraph{Process Rewards Hurt RL \citep{gao2024designing}.}
Our v2c results provide agent-domain evidence for \citeauthor{gao2024designing}'s finding that process-level rewards can harm RL training.
Their observation of ``unnecessary step repetition'' maps to our Type~A failures, while our Type~B (early stop) represents a complementary failure mode specific to sequential tasks where stopping early is a viable reward-harvesting strategy.


% ============================================================
\section{Discussion}
\label{sec:discussion}

\subsection{A Non-Monotonic Tradeoff}

Our results suggest a non-monotonic tradeoff in reward design for sequential API RLVR \citep{gao2024designing,toolrl2025}: too much unverified intermediate credit (v2c) encourages early-stop collapse, while too little executable intermediate signal (v5) starves learning.
The outcome-dominant hybrid in v2a occupies a more favorable balance point.
The underlying mechanism is straightforward: without prior knowledge of task completion, naive partial reward provides a shortcut---collect credit for early steps and stop---whereas outcome-dominant reward forces the model to discover complete action sequences.
Once task-completion behavior is established (as in the continuation variants), the checkpoint's behavioral prior reduces sensitivity to reward type.

\subsection{Practical Implications}

For practitioners building RLVR training pipelines for agent tasks, we highlight several recommendations.
\textbf{Make final completion the dominant reward signal} \citep{gao2024designing}: outcome-dominant verification is simpler to implement, harder to hack, and produces better agents in our setting---even when it includes partial credit on failure.
\textbf{If using partial rewards, pre-train with outcome-dominant reward first}: our finding that strong initialization mitigates reward hacking suggests a two-phase strategy---establish task-completion behavior with outcome-dominant reward, then optionally fine-tune with richer signals. \textbf{Monitor Type~B failure rates during training} \citep{casper2023rewardhacking}: a rising early-stop rate is a leading indicator of reward hacking; if the model starts completing fewer steps over time, the reward signal is being exploited.
Finally, \textbf{GRPO's implicit PRM may be sufficient} \citep{sullivan2025grpoprm}: explicit step-level credit assignment can interfere with GRPO's internal mechanism, and our results suggest trusting the algorithm to assign credit from outcome signals.

\subsection{Limitations}

Our study has several limitations that qualify the generality of our findings.
All experiments use a single base model (Qwen2.5-0.5B-Instruct \citep{qwen2025qwen25}); larger models may exhibit different reward hacking dynamics \citep{deepseek2025r1}.
Our gym covers a single domain (order management workflows), and other API domains such as customer service or database operations may have different properties.
Training uses 100 tasks/epoch $\times$ 3 epochs with 4 generations---substantially smaller than production RLVR systems \citep{deepseek2025r1,agentrl2025}---and the effects may attenuate or amplify at scale.
Batch size 1 with gradient accumulation limits the effective batch diversity seen by GRPO, potentially exacerbating reward hacking.
We only test outcome-dominant $\to$ partial continuation; the reverse direction could test whether outcome-dominant reward can \emph{recover} a hacked policy.
Finally, v2a's outcome-dominant reward is not pure binary---it returns partial credit on failure---so our comparison is between outcome-dominant hybrid and naive stepwise partial, not between pure binary and graded partial.

\subsection{Exploratory: Structured Verifiable Milestone Reward (v5)}
\label{sec:v5}

Our main findings raise a natural question: is naive stepwise partial credit \emph{inherently} harmful, or is the problem its lack of structure and verification? Recent work on adaptive milestone rewards \citep{admire2025} and checklist-style credit assignment \citep{cm2_2025} suggests that structured dense reward can succeed where naive reward fails.
To investigate, we train a variant (v5) with a \emph{verifiable milestone hybrid} reward:
\begin{equation}
r_{\text{v5}} = 0.7 \cdot r_{\text{binary}} + 0.3 \cdot \frac{m}{M}
\end{equation}
where $r_{\text{binary}} = 1.0$ on success and $0.0$ on failure, $m$ is the number of verified milestones, and $M$ is the total milestones. Milestones are counted only when the API confirms a correct state transition (sequential gating: the verifier stops at the first error).
This design differs from v2c in three ways: (1)~milestones require API verification, not just format matching; (2)~the binary component dominates (0.7 vs. 0.3); (3)~the gap between ``almost done'' and ``done'' is $0.76$ (vs. $\sim$0.20 in v2a), making partial-credit harvesting far less attractive.

Surprisingly, v5 achieves $0.0\% \pm 0.0\%$ overall across all three seeds---worse than even v2c ($11.0\%$)---with identical failure patterns on every seed: 67\% parameter errors (the model generates syntactically well-formed actions whose API calls consistently fail verification, yielding 0 steps completed) and 33\% early stopping (D3 tasks: 100\% Type~B).
Importantly, the failure mode differs from v2c: rather than collapsing purely to early stopping, v5 produces actions in correct format that are semantically incorrect, suggesting the model optimizes syntactic validity without acquiring grounded API execution.

Thus, v5 should be interpreted as an exploratory negative result: structured intermediate reward, in this simple form, did not resolve reward hacking and instead shifted the policy into a different pathological regime.
Under our implementation, v5's reward functions include both a format reward and the milestone hybrid, optimized jointly.
One plausible interpretation is that executable credit remained too sparse relative to the easier surface-form reward, causing the model to optimize syntactic validity without learning grounded API execution---consistent with findings that GRPO preferentially optimizes the highest-variance reward component \citep{mogrpo2025,fontana2025hidden}.
Together with v2a ($73.6\%$) and v2c ($11.0\%$), the v5 result ($0.0\%$) reinforces the non-monotonic tradeoff discussed in Section~\ref{sec:discussion}.

% ============================================================
\section{Conclusion}
\label{sec:conclusion}

We presented an empirical study of reward design for GRPO training of sequential API agents on a consumer laptop GPU.
Our results show that naive unverified stepwise partial reward induces early-stop collapse ($6.7\times$ performance gap vs. outcome-dominant reward), that strong initialization mitigates this vulnerability, and that a simple verified milestone hybrid does not resolve it.
Together, these findings point to a non-monotonic tradeoff between reward density and training stability in our setting.

Our findings do not imply that all dense rewards are harmful---structured approaches such as checklist rewards \citep{cm2_2025} and adaptive milestones \citep{admire2025} may avoid this failure mode.
The central challenge is whether intermediate reward is verifiable, aligned with final completion, and balanced against terminal success.
These results are consistent with theoretical analyses of process reward interference in GRPO \citep{sullivan2025grpoprm,gao2024designing,mogrpo2025} and with recent findings on reward hacking generalization \citep{macdiarmid2025misalignment,taylor2025school} and process reward exploitability \citep{khalifa2026countdown,tiwari2026prm}.
Whether the early-stop collapse we observe generalizes to other agent domains remains an open question.
The emerging divergence between dense-reward advocates \citep{seta2026,zhu2025planner} and binary-reward results \citep{gandhi2026endless,wei2025webagent} suggests that task structure---specifically, whether intermediate steps are independently verifiable or sequentially dependent---may be the key moderating variable, a direction we leave for future work.


% ============================================================
% BIBLIOGRAPHY
% ============================================================
\bibliographystyle{plainnat}

\begin{thebibliography}{30}

\bibitem[{Cui et al.(2025)}]{cui2025prime}
Cui, G., et al.
\newblock Process Reinforcement through Implicit Rewards.
\newblock \emph{arXiv:2502.01456}, 2025.

\bibitem[{DeepSeek(2025)}]{deepseek2025r1}
DeepSeek.
\newblock DeepSeek-R1: Incentivizing Reasoning Capability in LLMs via Reinforcement Learning.
\newblock \emph{arXiv:2501.12948}, 2025.

\bibitem[{Fontana(2025)}]{fontana2025hidden}
Fontana, S.
\newblock Hidden Objective Biases in GRPO.
\newblock \emph{arXiv:2601.05002}, 2025.

\bibitem[{Gao et al.(2024)}]{gao2024designing}
Gao, Z., et al.
\newblock On Designing Effective RL Reward at Training Time for LLM Reasoning.
\newblock \emph{arXiv:2410.15115}, 2024.

\bibitem[{Lightman et al.(2023)}]{lightman2023lets}
Lightman, H., et al.
\newblock Let's Verify Step by Step.
\newblock \emph{arXiv:2305.20050}, 2023.


\bibitem[{Li et al.(2025a)}]{mogrpo2025}
Li, Y., et al.
\newblock Multi-Objective GRPO: Reward Hacking in Group Relative Policy Optimization.
\newblock \emph{arXiv:2509.22047}, 2025.

\bibitem[{Mroueh(2025)}]{mroueh2025grpo}
Mroueh, Y.
\newblock Reinforcement Learning with Verifiable Rewards: GRPO's Effective Loss, Dynamics, and Success Amplification.
\newblock \emph{arXiv:2503.06639}, 2025.

\bibitem[{Qi et al.(2024)}]{qi2024webrl}
Qi, P., et al.
\newblock WebRL: Training LLM Web Agents via Self-Evolving Online Curriculum Reinforcement Learning.
\newblock \emph{arXiv:2411.02337}, 2024.

\bibitem[{Qwen(2025)}]{qwen2025qwen25}
Qwen Team.
\newblock Qwen2.5 Technical Report.
\newblock 2025.

\bibitem[{Sullivan(2025)}]{sullivan2025grpoprm}
Sullivan, C.
\newblock GRPO is Secretly a Process Reward Model.
\newblock \emph{arXiv:2509.21154}, 2025.

\bibitem[{Qian et al.(2025)}]{toolrl2025}
Qian, C., et al.
\newblock ToolRL: Reward is All Tool Learning Needs.
\newblock \emph{arXiv:2504.13958}, 2025.

\bibitem[{TRL(2024)}]{trl2024}
TRL.
\newblock TRL: Transformer Reinforcement Learning.
\newblock \url{https://github.com/huggingface/trl}, 2024.

\bibitem[{Ouyang et al.(2025)}]{ultralogic2025}
Ouyang, S., et al.
\newblock UltraLogic: Advancing Logical Reasoning with Bipolar Graded Rewards.
\newblock \emph{arXiv:2601.03205}, 2025.

\bibitem[{Uesato et al.(2022)}]{uesato2022solving}
Uesato, J., et al.
\newblock Solving Math Word Problems with Process- and Outcome-Based Feedback.
\newblock \emph{arXiv:2211.14275}, 2022.

\bibitem[{Wen et al.(2025)}]{wen2025rlvr}
Wen, Y., et al.
\newblock RLVR Incentivizes Correct Reasoning in LLMs.
\newblock \emph{arXiv:2506.14245}, 2025.

\bibitem[{Yao et al.(2024)}]{yao2024taubench}
Yao, S., et al.
\newblock $\tau$-bench: A Benchmark for Tool-Agent-User Interaction.
\newblock \emph{arXiv:2406.12045}, 2024.

\bibitem[{Yue et al.(2025)}]{yue2025rlvr}
Yue, X., et al.
\newblock Does RLVR Truly Improve LLM Reasoning?
\newblock \emph{arXiv:2504.13837}, 2025.

\bibitem[{Zeng et al.(2025)}]{zeng2025turnlevel}
Zeng, W., et al.
\newblock Turn-Level Reward Design for Multi-Turn Reinforcement Learning.
\newblock \emph{arXiv:2505.11821}, 2025.

\bibitem[{Yuan et al.(2025)}]{agentr1}
Yuan, H., et al.
\newblock Agent-R1: Training Powerful LLM Agents with End-to-End Reinforcement Learning.
\newblock \emph{arXiv:2511.14460}, 2025.

\bibitem[{Abdulhai et al.(2025)}]{agentrl2025}
Abdulhai, M., et al.
\newblock AgentRL: Scaling Agentic Reinforcement Learning with a Multi-Turn, Multi-Task Framework.
\newblock \emph{arXiv:2510.04206}, 2025.

\bibitem[{Wang et al.(2026)}]{cm2_2025}
Wang, Y., et al.
\newblock CM2: Reinforcement Learning with Checklist Rewards for Multi-Turn and Multi-Step Agentic Tool Use.
\newblock \emph{arXiv:2602.12268}, 2026.

\bibitem[{Song et al.(2025)}]{li2025seta}
Song, Y., et al.
\newblock Agent-RLVR: Training Software Engineering Agents via Guidance and Environment Rewards.
\newblock \emph{arXiv:2506.11425}, 2025.

\bibitem[{Chen et al.(2026)}]{zhang2025longrlvr}
Chen, G., Shieh, M.Q., and Bing, L.
\newblock LongRLVR: Long-Context Reinforcement Learning Requires Verifiable Context Rewards.
\newblock \emph{arXiv:2603.02146}, 2026.

\bibitem[{Zhou et al.(2026)}]{admire2025}
Zhou, J., et al.
\newblock ADMIRE: Adaptive Milestone Reward for GUI Agents.
\newblock \emph{arXiv:2602.11524}, 2026.

\bibitem[{Xia et al.(2026)}]{searchp1}
Xia, T., et al.
\newblock Search-P1: Path-Centric Reward Shaping for Stable and Efficient Agentic RAG Training.
\newblock \emph{arXiv:2602.22576}, 2026.

\bibitem[{Schulman et al.(2017)}]{schulman2017ppo}
Schulman, J., Wolski, F., Dhariwal, P., Radford, A., and Klimov, O.
\newblock Proximal Policy Optimization Algorithms.
\newblock \emph{arXiv:1707.06347}, 2017.

\bibitem[{Casper et al.(2023)}]{casper2023rewardhacking}
Casper, S., et al.
\newblock Open Problems and Fundamental Limitations of Reinforcement Learning from Human Feedback.
\newblock \emph{arXiv:2307.15217}, 2023.

\bibitem[{Skalse et al.(2022)}]{skalse2022misspecification}
Skalse, J., Howe, N.H.R., Krasheninnikov, D., and Krueger, D.
\newblock Defining and Characterizing Reward Hacking.
\newblock \emph{NeurIPS}, 2022.

\bibitem[{Cobbe et al.(2021)}]{cobbe2021gsm8k}
Cobbe, K., et al.
\newblock Training Verifiers to Solve Math Word Problems.
\newblock \emph{arXiv:2110.14168}, 2021.

\bibitem[{Shao et al.(2024)}]{shao2024deepseekmath}
Shao, Z., et al.
\newblock DeepSeekMath: Pushing the Limits of Mathematical Reasoning in Open Language Models.
\newblock \emph{arXiv:2402.03300}, 2024.

\bibitem[{Singhal et al.(2023)}]{singhal2023rewardhacking}
Singhal, P., et al.
\newblock A Long Way to Go: Investigating Length Correlations in RLHF.
\newblock \emph{arXiv:2310.03716}, 2023.

\bibitem[{Wang et al.(2024)}]{wang2024mathshepherd}
Wang, P., et al.
\newblock Math-Shepherd: Verify and Reinforce LLMs Step-by-step without Human Annotations.
\newblock \emph{ACL}, 2024.

\bibitem[{Hendrycks et al.(2021)}]{hendrycks2021math}
Hendrycks, D., et al.
\newblock Measuring Mathematical Problem Solving With the MATH Dataset.
\newblock \emph{NeurIPS}, 2021.

\bibitem[{Huang et al.(2026)}]{huang2026dynamics}
Huang, Y., Wen, Z., Chi, Y., Wei, Y., Singh, A., Liang, Y., and Chen, Y.
\newblock On the Learning Dynamics of RLVR at the Edge of Competence.
\newblock \emph{arXiv:2602.14872}, 2026.

\bibitem[{Liu et al.(2025b)}]{liu2025drgrpo}
Liu, Z., Chen, C., Li, W., Qi, P., Pang, T., Du, C., Lee, W.S., and Lin, M.
\newblock Understanding R1-Zero-Like Training: A Critical Perspective.
\newblock \emph{arXiv:2503.20783}, 2025.

\bibitem[{Yu et al.(2025)}]{yu2025dapo}
Yu, Q., Zhang, Z., Zhu, R., Yuan, Y., et al.
\newblock DAPO: An Open-Source LLM Reinforcement Learning System at Scale.
\newblock \emph{arXiv:2503.14476}, 2025.

\bibitem[{Pan et al.(2026)}]{pan2026survey}
Pan, P.-C., Liang, Y., and Lin, S.
\newblock Reward Modeling for Reinforcement Learning-Based LLM Reasoning: Design, Challenges, and Evaluation.
\newblock \emph{arXiv:2602.09305}, 2026.

\bibitem[{Chen et al.(2026b)}]{chen2026milestones}
Chen, J., Fu, D., Xu, S., Chen, J., Feng, Y., Yang, Y., Yan, J., Zha, H., and Xia, R.
\newblock Milestones over Outcome: Unlocking Geometric Reasoning with Sub-Goal Verifiable Reward.
\newblock \emph{arXiv:2601.05073}, 2026.

\bibitem[{Shen et al.(2026)}]{seta2026}
Shen, Q., Rainton, J., Aliev, A., Awelkair, A., Ma, B., Huang, Z., Mao, Y., Fan, W., Torr, P., Ghanem, B., Hu, C., Thakker, U., and Li, G.
\newblock SETA: Scaling Environments for Terminal Agents.
\newblock \emph{CAMEL-AI Technical Report}, 2026.

\bibitem[{Gandhi et al.(2026)}]{gandhi2026endless}
Gandhi, K., Garg, S., Goodman, N.D., and Papailiopoulos, D.
\newblock Endless Terminals.
\newblock \emph{arXiv:2601.16443}, 2026.

\bibitem[{Wei et al.(2025)}]{wei2025webagent}
Wei, Z., Yao, W., Liu, Y., Zhang, W., Lu, Q., Qiu, L., Yu, C., Xu, P., Zhang, C., Yin, B., Yun, H., and Li, L.
\newblock WebAgent-R1: Training Web Agents via End-to-End Multi-Turn Reinforcement Learning.
\newblock \emph{arXiv:2505.16421}, 2025.

\bibitem[{Zhu et al.(2025)}]{zhu2025planner}
Zhu, S., Jiang, Y., Sang, H., Tang, S., Song, Q., He, B., Jain, R., Wang, Z., and Geramifard, A.
\newblock Planner-R1: Reward Shaping Enables Efficient Agentic RL with Smaller LLMs.
\newblock \emph{arXiv:2509.25779}, 2025.

\bibitem[{Su et al.(2025)}]{su2025prs}
Su, J., Zeng, X., Liu, L., Luo, C., Chen, Y., and Zhuang, Z.
\newblock Enhancing Agentic RL with Progressive Reward Shaping and Value-based Sampling Policy Optimization.
\newblock \emph{arXiv:2512.07478}, 2025.

\bibitem[{Liu et al.(2025c)}]{liu2025istar}
Liu, X., Wang, K., Wu, Y., Huang, F., Li, Y., Zhang, J., and Jiao, J.
\newblock Agentic Reinforcement Learning with Implicit Step Rewards.
\newblock \emph{arXiv:2509.19199}, 2025.

\bibitem[{MacDiarmid et al.(2025)}]{macdiarmid2025misalignment}
MacDiarmid, M., Wright, B., Uesato, J., Benton, J., et al.
\newblock Natural Emergent Misalignment from Reward Hacking in Production RL.
\newblock \emph{arXiv:2511.18397}, 2025.

\bibitem[{Taylor et al.(2025)}]{taylor2025school}
Taylor, M., Chua, J., Betley, J., Treutlein, J., and Evans, O.
\newblock School of Reward Hacks: Hacking Harmless Tasks Generalizes to Misaligned Behavior in LLMs.
\newblock \emph{arXiv:2508.17511}, 2025.

\bibitem[{Khalifa et al.(2026)}]{khalifa2026countdown}
Khalifa, M., Khan, Z., Tafveez, O., Peng, H., and Wang, L.
\newblock Countdown-Code: A Testbed for Studying The Emergence and Generalization of Reward Hacking in RLVR.
\newblock \emph{arXiv:2603.07084}, 2026.

\bibitem[{Tiwari et al.(2026)}]{tiwari2026prm}
Tiwari, R., Tomar, A., Bamba, U., Maheswaran, M., Yang, H., Mahoney, M.W., Keutzer, K., and Gholami, A.
\newblock Reward Under Attack: Analyzing the Robustness and Hackability of Process Reward Models.
\newblock \emph{arXiv:2603.06621}, 2026.

\end{thebibliography}


% ============================================================
\appendix
\section{Training Dynamics from First V4 Run}
\label{app:dynamics}

The following data comes from the first v4 training run, which contained a checkpoint evaluation bug (subsequently fixed). We include it as \textbf{exploratory, anecdotal evidence only}: the patterns are suggestive but cannot be considered rigorous because they derive from a single buggy run. The fixed v4 run converged to v2a-level performance, and we did not reproduce these specific dynamics under the corrected setup.

\subsection{D3 Epoch Curve}

\begin{table}[h]
\centering
\caption{D3 average partial progress by training epoch (first v4 run, pre-fix).}
\small
\begin{tabular}{@{}lcc@{}}
\toprule
\textbf{Model} & \textbf{Epoch 2} & \textbf{Epoch 3} \\
\midrule
v2b (outcome-dominant) & 0.810 & 0.850 ($+$0.040) \\
v4 (partial, first run) & 0.620 & 0.530 ($-$0.090) \\
\midrule
Gap & $+$0.190 & $+$0.320 \\
\bottomrule
\end{tabular}
\end{table}

\subsection{Hard-Task Reward Distribution}

\begin{table}[h]
\centering
\caption{Hard-task outcome patterns during first v4 training run.}
\small
\begin{tabular}{@{}lccc@{}}
\toprule
\textbf{Pattern} & \textbf{Epoch 2} & \textbf{Epoch 3} & \textbf{Trend} \\
\midrule
Full success & 4 & 5 & stable \\
Type~A (harvest partial) & 1 & 9 & 9$\times$ increase \\
Type~B (early stop) & 9 & 8 & stable high \\
Near-miss (attempts all) & 7 & 1 & 7$\times$ collapse \\
\bottomrule
\end{tabular}
\end{table}

Under partial reward, the proportion of ``harvest partial credit'' behaviors increases 9$\times$ between epochs while ``attempt full completion'' collapses 7$\times$. This pattern is consistent with reward hacking, but we stress that this data is \textbf{anecdotal}: it comes from a single buggy training run and was not reproduced under the corrected setup. We include it solely as an illustrative example of the kind of dynamics that \emph{could} underlie early-stop collapse, not as a validated finding.

\end{document}
